% Generated from Review/corrected_text.md (the reviewed source text).

\bakhsh{پیوست الف - سیستم ترکیبی برایتول}

\band{۱. کلیات}

سیستم ترکیبی برایتول روشی برای رفع تساوی در مسابقات اتللو است که با سیستم سوئیسی برگزار می‌شوند. این روش در مسابقات دوره‌ای نیز قابل استفاده است و در آنجا معادل شمارش مهره‌هاست. در اصل، این روش ترکیبی از قدرت حریفان و تعداد مهره‌هاست.

\band{۲. ثابت \lr{C}}

سیستم \lr{MBS} به ثابت \lr{C} بستگی دارد که وزن نسبی اهمیت قدرت حریفان و تعداد مهره‌ها را تعیین می‌کند. در مسابقات جهانی اتللو، با فرض اینکه مسابقات به شکل دوره‌ای یا ۱۳ دور سیستم سوئیسی برگزار شود، مقدار \lr{C} برابر ۶ است.

\band{۳. ضریب برایتول}

در پایان مسابقات، مسئولان برگزاری برای هر بازیکنی که در تساوی قرار دارد، ضریب برایتول را محاسبه می‌کنند. نحوه محاسبه به شرح زیر است:

\begin{itemize}
  \item اگر یکی از حریفان بازیکن در طول مسابقات انصراف داده است، بازی با او را فعلاً در نظر نگیرید. همچنین اگر بازیکن با \lr{BYE} جفت شده (یا یک دور را بازی نکرده) است، آن بازی را در نظر نگیرید.
  \item مجموع مهره‌های بازیکن را در تمام بازی‌هایی که در مرحله قبل کنار گذاشته نشده‌اند حساب کنید و حاصل‌ضرب \lr{C} در مجموع امتیازاتی را که تمام حریفان بازیکن (به جز حریفانی که انصراف داده‌اند) در مسابقات کسب کرده‌اند، به آن اضافه کنید.
  \item به ازای هر بازی در برابر حریفی که انصراف داده است و هر استراحت (\lr{bye}) یا دور بازی‌نشده، عدد ۳۲ به علاوه (\lr{C} ضرب در امتیاز خود بازیکن در مسابقات) را به نتیجه مرحله قبل اضافه کنید. عدد حاصل، ضریب برایتول بازیکن است.
\end{itemize}

\band{۴. رفع تساوی امتیازات}

اگر دو یا چند بازیکن در مسابقات امتیاز برابر داشته باشند، تساوی با رده‌بندی بازیکنان به ترتیب نزولی ضریب برایتول رفع می‌شود. اگر دو یا چند بازیکن هم امتیاز برابر و هم ضریب برایتول برابر داشته باشند، رده‌بندی به ترتیب نزولی تعداد مهره‌ها انجام خواهد شد. در صورت ادامه تساوی، رده‌بندی با یک روش تصادفی عادلانه (پرتاب یک مهره اتللو یا بیرون کشیدن ژتون‌های شماره‌دار از یک کیسه دربسته) مشخص می‌شود.
